\documentclass[11pt]{article}
\usepackage[margin=1in]{geometry}
\usepackage{amsmath,amssymb,graphicx,booktabs,hyperref,natbib}
\title{Dynamic Multi-Asset Portfolio Optimization Across the Accumulation--Decumulation Lifecycle:\\
A Utility-Based Analysis of Savings and Withdrawal Rates}
\author{Tanish Patel\\Columbia University \and Wedmont Private Capital}
\date{Draft --- figures built from stored run IDs (\texttt{make paper})}

\begin{document}
\maketitle

\begin{abstract}
We measure how much lifetime utility an investor gains from a cash-flow-aware, dynamically optimised multi-asset
portfolio relative to a static 60/40 allocation or an age-based glidepath, expressed as the percentage gain in
certainty-equivalent (CE) real consumption. Policies are compared on common random numbers; Monte Carlo, parameter
and model uncertainty are reported separately, and a result is called robust only if its sign survives all three.
\end{abstract}

\section{Introduction}
Most retirement tools ask whether a portfolio survives. We ask how lifetime welfare changes when the allocation
responds to the investor's own cash flows, horizon and economic regime \citep{merton1969lifetime,cocco2005consumption}.

\section{Data}
Investable returns come from WRDS (CRSP index and Treasury files, WRDS Bond Returns, ETFs resolved by PERMNO);
macro state variables, yield curves \citep{gurkaynak2007us,gurkaynak2010tips} and term premia \citep{adrian2013pricing}
from public sources. Two universes: a long-history core and a modern diversified universe.

\section{Model}
Wealth evolves monthly in real terms,
\begin{equation}
W_{t+1}=(W_t+Y_t+P_t-C_t)(1+\mathbf w_t^\top\mathbf R_{t+1})-\kappa\lVert\mathbf w_t-\tilde{\mathbf w}_t\rVert_1(W_t+Y_t+P_t-C_t),
\end{equation}
and the investor maximises survival-weighted CRRA utility with a bequest motive. The headline metric is
\begin{equation}
\mathrm{CE}(\pi)=\Big[\tfrac{(1-\gamma)\hat U(\pi)}{\sum_t\beta^tS_t}\Big]^{1/(1-\gamma)},\qquad
\Delta_{\mathrm{CE}}=\mathrm{CE}(\pi^*)/\mathrm{CE}(\pi_{\mathrm{bench}})-1.
\end{equation}
Scenarios come from a stationary block bootstrap \citep{politis1994stationary,politis2004automatic}, a multivariate
Student-t, a Markov-switching model \citep{hamilton1989new} and a yield-curve factor model; dependence is studied with
DCC-GARCH \citep{engle2002dynamic} and Ledoit--Wolf shrinkage \citep{ledoit2004well}.

\section{Hypotheses}
\begin{description}
\item[H1] Optimal allocation depends on the sign and magnitude of the net-flow rate $q_t$, not only age.
\item[H2] A lifecycle-optimised policy raises CE consumption over constant-mix and target-date policies.
\item[H3] The CE value of each sleeve differs between inflationary and disinflationary regimes.
\item[H4] Sleeves beyond US equity and Treasuries improve consumption sustainability under some conditions.
\end{description}

\section{Results}
\begin{figure}[h]\centering
\IfFileExists{figures/fig1_decumulation_heatmap.pdf}{\includegraphics[width=.9\linewidth]{figures/fig1_decumulation_heatmap.pdf}}{\fbox{run \texttt{make paper}}}
\caption{$\Delta_{\mathrm{CE}}$ of the parametric dynamic policy vs 60/40: withdrawal rate $\times$ years relative to retirement.}
\end{figure}
\begin{figure}[h]\centering
\IfFileExists{figures/fig2_accumulation_heatmap.pdf}{\includegraphics[width=.9\linewidth]{figures/fig2_accumulation_heatmap.pdf}}{\fbox{run \texttt{make paper}}}
\caption{Accumulation: savings rate $\times$ starting age, vs the target-date glidepath.}
\end{figure}
\begin{figure}[h]\centering
\IfFileExists{figures/fig3_transition.pdf}{\includegraphics[width=.75\linewidth]{figures/fig3_transition.pdf}}{\fbox{run \texttt{make paper}}}
\caption{Lifecycle transition: optimal equity weight against $q_t$ (grid DP, reduced problem).}
\end{figure}
\begin{figure}[h]\centering
\IfFileExists{figures/fig4_shapley_by_regime.pdf}{\includegraphics[width=.75\linewidth]{figures/fig4_shapley_by_regime.pdf}}{\fbox{run \texttt{make seed-full}}}
\caption{Shapley marginal utility of sleeves by starting regime.}
\end{figure}

\section{Robustness}
Monte Carlo convergence, parameter bootstrap re-estimation, model disagreement across generators, calibration
windows and historical cohort replay \citep{bengen1994determining,guyton2006decision}.

\bibliographystyle{plainnat}
\bibliography{references}
\end{document}
